\chapter*{\xiaosan\heiti{摘\quad 要}}
\addcontentsline{toc}{chapter}{摘要}

路网是城市交通系统的核心基础设施，其表示学习对于路径规划、行程时间估计和交通速度推断等智能交通任务具有重要意义。真实路网呈现出显著的层次结构特性：从个体路段到结构区域，再到功能地理区域，形成自然的树状层级组织，并伴有幂律度分布。然而，现有的基于欧氏空间的路网表示学习方法在编码此类层次结构时存在本质局限——欧氏空间的多项式体积增长与路网层次结构的指数分支模式之间存在根本性不匹配，导致嵌入失真严重，层次关系难以精确建模。

双曲几何以其负曲率和指数体积增长特性，为层次结构数据的嵌入提供了天然的几何先验。本文以双曲几何为统一工具，围绕路网层次结构建模展开研究：从利用双曲空间几何特性隐式捕获层次信息，到通过蕴含锥机制显式约束层次关系，逐步深化对路网层次结构的几何建模能力。

本文首先对路网层次结构特性与双曲空间的适配性进行了系统的理论分析。通过对29个国家50个城市的大规模 $\delta$-双曲性分析，实证验证了路网拓扑结构的树状特性，并从指数体积增长、层次深度编码和父子关系建模三个角度论证双曲空间作为路网层次建模几何框架的理论优势。在此基础上，提出了层次建模的两种范式——隐式捕获与显式约束，为后续两项工作提供了统一的理论根基。

第一项工作提出了双曲路网图神经网络（HRGNN），采用隐式层次建模范式。HRGNN 利用庞加莱球模型建模路网的层次结构，通过双通道双曲注意力机制同时捕获局部邻域的连通模式和重要性引导的深度优先搜索（DFS）路径所揭示的层次依赖关系。自适应融合策略根据路网特征动态调整两条路径的贡献权重。在多个真实城市路网的道路类型分类实验中，HRGNN 在城市内分类任务上相比欧氏基线提升微平均 F1 分数 8.4\%--9.1\%，在跨城市迁移学习任务上相比现有双曲方法提升宏平均 F1 分数 13.0\%，验证了双曲几何对路网层次结构建模的有效性。

第二项工作提出了基于蕴含锥的双曲层次路网表示模型（HHRoad），采用显式层次建模范式，进一步解决隐式建模中层次关系缺乏几何约束的不足。HHRoad 在 Lorentz 双曲模型中运行，构建路段（Segment）、局部区域（Locality）和功能区（Region）三层层次结构，通过三项核心创新实现显式的层次约束：（1）双曲嵌入层将路段属性映射至 Lorentz 超球面，以到原点的距离自然编码层次深度；（2）组合蕴含学习，利用蕴含锥机制对 Region$\to$Locality$\to$Segment 的层次关系施加显式约束；（3）基于 Lorentz 距离的层次对比学习，同时捕获层内和跨层的语义相似性。在 VecCity 基准框架下对五个真实城市数据集的实验表明，HHRoad 在八种方法中取得最佳综合平均排名（3.07），在轨迹相似性搜索任务上相比最强基线的 ACC@3 最高提升 9.4\%。

本文的两项工作共同表明：双曲几何能够为路网的层次结构建模提供有效的几何框架，从隐式利用空间特性到显式设计几何约束，能够适配路网天然的树状拓扑，以更低的失真编码层次关系，并带来跨任务、跨城市的性能提升。

\vspace{0.5cm}
\sihao{\heiti{关键词：}}\xiaosi{路网表示学习；双曲几何；图神经网络；层次结构建模；蕴含锥；对比学习}